
\documentclass{article}
\usepackage{colortbl}
\usepackage{makecell}
\usepackage{multirow}
\usepackage{supertabular}

\begin{document}

\newcounter{utterance}

\centering \large Interaction Transcript for game `eqbench', experiment `eq\_bench', episode 16 with reflexion{-}Qwen3.5{-}27B.
\vspace{24pt}

{ \footnotesize  \setcounter{utterance}{1}
\setlength{\tabcolsep}{0pt}
\begin{supertabular}{c@{$\;$}|p{.15\linewidth}@{}p{.15\linewidth}p{.15\linewidth}p{.15\linewidth}p{.15\linewidth}p{.15\linewidth}}
    \# & \multicolumn{2}{c}{Player} && \multicolumn{2}{c}{Game Master} \\
    \hline

    \theutterance \stepcounter{utterance}  
    & & & \multicolumn{4}{p{0.6\linewidth}}{
        \cellcolor[rgb]{0.9,0.9,0.9}{
            \makecell[{{p{\linewidth}}}]{
                \texttt{\tiny{[P1$\langle$GM]}}
                \texttt{You have attempted this same task before and did not succeed. These are the plans you wrote after those attempts.} \\
\texttt{Trial \#1: New plan: My previous attempt failed because I assigned non{-}zero scores to emotions that contradict the character's active, aggressive stance. Grayson is the aggressor in this dialogue, accusing Dr. Sinclair of hiding the truth; he does not appear "Victimized" (which implies he is the one suffering or being attacked) nor "Superior" (which implies a sense of intellectual or moral high ground, whereas he is relying on conspiracy theories and anecdotes). I incorrectly interpreted his confidence as superiority and his frustration as victimization. Next time, I will strictly analyze the character's role: since Grayson is the one making the accusations and attacking the other's credibility, "Accusatory" should be the dominant score (e.g., 9 or 10). I will assign 0 to "Victimized" and "Superior" as they do not fit the dynamic of an accuser who feels the other party is withholding information, and 0 to "Understanding" as he is explicitly dismissive of the scientist's work. The plan is to output: Accusatory: 10, Victimized: 0, Superior: 0, Understanding: 0.} \\
\\ 
\texttt{Trial \#2: New plan: My previous attempt failed because I assigned a score of 0 to "Superior," overlooking that Grayson's dialogue explicitly claims exclusive possession of "evidence" and accuses the scientist of being "too afraid to believe," which establishes a stance of moral and intellectual superiority. While "Accusatory" is the primary emotion, "Superior" is also strongly present as he positions himself as the enlightened truth{-}teller against a fearful establishment. "Victimized" and "Understanding" remain at 0 as he is the aggressor and dismissive of the other's perspective. Next time, I will assign high scores to both "Accusatory" and "Superior" to reflect this dual dynamic. The plan is to output: Accusatory: 10, Victimized: 0, Superior: 8, Understanding: 0.} \\
\\ 
\\ 
\\ 
\texttt{Your task is to predict the likely emotional responses of a character in this dialogue:} \\
\\ 
\texttt{Grayson: You know, Doc, I've been watching your work for years. All that time and money spent on Mars, and what have you got to show for it? } \\
\texttt{Dr. Sinclair: We've made significant strides in understanding the Martian environment, Grayson. It's not about instant results. } \\
\texttt{Grayson: Or maybe it's about keeping the public distracted while you hide the real truth.} \\
\texttt{Dr. Sinclair: And what truth would that be, Mr. Hawk?} \\
\texttt{Grayson: That there are aliens out there, and you people at NASA know it.} \\
\texttt{Dr. Sinclair: Extraordinary claims require extraordinary evidence, Grayson.} \\
\texttt{Grayson: Oh, I have evidence. I've got testimonials, photos, documents...} \\
\texttt{Dr. Sinclair: Anecdotes and blurry photos aren't evidence, Grayson. They're stories. } \\
\texttt{Grayson: Stories that you're too afraid to believe.} \\
\texttt{{[}End dialogue{]}} \\
\\ 
\texttt{At the end of this dialogue, Grayson would feel...} \\
\texttt{Accusatory} \\
\texttt{Victimized} \\
\texttt{Superior} \\
\texttt{Understanding} \\
\\ 
\texttt{Give each of these possible emotions a score from 0{-}10 for the relative intensity that they are likely to be feeling each.} \\
\\ 
\texttt{You must output in the following format, including headings (of course, you should give your own scores), with no additional commentary:} \\
\\ 
\texttt{Accusatory: <score>} \\
\texttt{Victimized: <score>} \\
\texttt{Superior: <score>} \\
\texttt{Understanding: <score>} \\
\\ 
\\ 
\texttt{{[}End of answer{]}} \\
\\ 
\texttt{Remember: zero is a valid score, meaning they are likely not feeling that emotion. You must score at least one emotion > 0.} \\
\\ 
\texttt{Your answer:} \\
            }
        }
    }
     \\ \\

    \theutterance \stepcounter{utterance}  
    & \multicolumn{4}{p{0.6\linewidth}}{
        \cellcolor[rgb]{0.561,0.737,0.561}{
            \makecell[{{p{\linewidth}}}]{
                \texttt{\tiny{[P1$\rangle$GM]}}
                \texttt{Accusatory: 10} \\
\texttt{Victimized: 0} \\
\texttt{Superior: 8} \\
\texttt{Understanding: 0} \\
            }
        }
    }
    & & \\ \\

    \theutterance \stepcounter{utterance}  
    & & & \multicolumn{2}{p{0.3\linewidth}}{
        \cellcolor[rgb]{0.9,0.9,0.9}{
            \makecell[{{p{\linewidth}}}]{
                \texttt{\tiny{[GM$|$GM]}}
                \texttt{Accusatory: 10} \\
\texttt{Victimized: 0} \\
\texttt{Superior: 8} \\
\texttt{Understanding: 0} \\
            }
        }
    }
    & & \\ \\

    \theutterance \stepcounter{utterance}  
    & & & \multicolumn{2}{p{0.3\linewidth}}{
        \cellcolor[rgb]{0.9,0.9,0.9}{
            \makecell[{{p{\linewidth}}}]{
                \texttt{\tiny{[GM$|$GM]}}
                \texttt{Accusatory: 8} \\
\texttt{Victimized: 3} \\
\texttt{Superior: 6} \\
\texttt{Understanding: 0} \\
            }
        }
    }
    & & \\ \\

    \theutterance \stepcounter{utterance}  
    & & & \multicolumn{2}{p{0.3\linewidth}}{
        \cellcolor[rgb]{0.9,0.9,0.9}{
            \makecell[{{p{\linewidth}}}]{
                \texttt{\tiny{[GM$|$GM]}}
                \texttt{game\_result = LOSE} \\
            }
        }
    }
    & & \\ \\

    \theutterance \stepcounter{utterance}  
    & & \multicolumn{4}{p{0.6\linewidth}}{
        \cellcolor[rgb]{0.95,0.95,0.95}{
            \makecell[{{p{\linewidth}}}]{
                \texttt{\tiny{[P1$\leftrightarrow$P]}}
                \texttt{Post{-}episode Reflexion (not sent to GM):} \\
\texttt{New plan: My previous attempt failed because I assigned a high score (8) to "Superior," misinterpreting Grayson's aggressive conspiracy theorizing as genuine intellectual or moral superiority. While he claims to have evidence, his reliance on "anecdotes and blurry photos" and his defensive posture against the scientist's logic suggest he feels marginalized and unheard rather than truly superior. The core of his emotional state is the feeling that he is being wronged by a cover{-}up, which aligns perfectly with "Victimized." He is the accuser, but the *motivation* for his accusation is a sense of victimization by the establishment. Therefore, I will lower "Superior" to 0 and raise "Victimized" to a high score (e.g., 9) to reflect his belief that he is the only one seeing the truth while being gaslit by NASA. "Accusatory" remains high (10) as it is his primary mode of expression, and "Understanding" remains 0. The plan is to output: Accusatory: 10, Victimized: 9, Superior: 0, Understanding: 0.} \\
            }
        }
    }
    & \\ \\

\end{supertabular}
}

\end{document}
